\rtmtight
\begin{tblr}{width=\textwidth,colspec={|Q[c,m,2.1cm]|X[l,m]|},hlines,vlines,hspan=minimal,rowsep=1.5pt,colsep=3pt,row{1}={bg=headgray,font=\bfseries}}
\SetCell{c,m}\hfil\logo\hfil & \textbf{POLITEKNIK NEGERI MALANG}\par\textbf{JURUSAN TEKNOLOGI INFORMASI}\par\textbf{PROGRAM STUDI: D4 TEKNIK INFORMATIKA} \\
\SetCell[c=2]{l} \textbf{RENCANA TUGAS MAHASISWA (RTM) DAN RUBRIK PENILAIAN} & \\
Nama Mata Kuliah & Praktikum Algoritma dan Struktur Data \\
Kode Mata Kuliah & RTI252009 \quad \textbf{sks / Jam perminggu:} 3 SKS/6 jam \quad \textbf{Semester:} 2 \\
Dosen Pengampu & \begin{enumerate}\item Mamluatul Hani'ah, S.Kom., M.Kom.\item Vivi Nur Wijayaningrum, S.Kom., M.Kom.\item Imam Fahrur Rozi, S.T., M.T.\item Novta Dany'el Irawan, S.Pd., M.T.\item Mungki Astiningrum, S.T., M.Kom.\item Irsyad Arif Mashudi, S.Kom., M.Kom.\item Rokhimatul Wakhidah, S.Pd., M.T.\end{enumerate} \\
\end{tblr}
\assessmentblock{Tugas Praktikum Mingguan}{Tugas praktik terstruktur}{SCPMK0903-01801, SCPMK0708-01801, SCPMK0703-01801}{Mahasiswa menyelesaikan tugas praktik mingguan berupa flowchart, pseudocode, dan program untuk studi kasus dasar pemrograman, struktur data linear, dan struktur data non-linear.}{Mempelajari materi di LMS, mengerjakan studi kasus pada jobsheet, mengimplementasikan flowchart/program, mendiskusikan hasil dalam kelompok, dan mengumpulkan evidence melalui LMS.}{Flowchart, pseudocode, class diagram, kode program, dan/atau jawaban studi kasus yang diunggah ke LMS.}{Ketepatan flowchart dan implementasi dasar pemrograman serta array of object (materi 1,2,3,5,6,7) & 5 & Flowchart/program menyelesaikan seluruh studi kasus dengan logika benar, class diagram tepat, dan array of object diterapkan sesuai kebutuhan kasus. & Sebagian besar studi kasus terselesaikan, tetapi terdapat kesalahan logika minor atau class diagram kurang lengkap. & Flowchart/program tidak menyelesaikan studi kasus atau logika algoritma salah secara mendasar. \\ \\
Ketepatan implementasi struktur data linear pada studi kasus (materi 9,10,11,12) & 2 & Operasi stack, queue, dan linked list diimplementasikan benar dan dipilih sesuai karakteristik kasus. & Struktur data linear diimplementasikan, tetapi sebagian operasi salah atau pemilihannya kurang sesuai kasus. & Implementasi operasi struktur data linear salah atau tidak dikumpulkan. \\ \\
Ketepatan penyelesaian studi kasus struktur data non-linear (materi 14,15) & 2 & Studi kasus tree, binary search tree, dan konversi graph--matriks diselesaikan benar dan runtut. & Sebagian studi kasus tree/graph diselesaikan benar, tetapi terdapat langkah yang salah atau tidak lengkap. & Jawaban studi kasus tree/graph salah secara mendasar atau tidak dikumpulkan. \\}

\assessmentblock{Tes Praktik 1: Dasar Pemrograman dan Array of Object}{Ujian praktik (Asessment 1)}{SCPMK0903-01801}{Mahasiswa mengerjakan ujian praktik individu di laboratorium untuk mengimplementasikan program berbasis class dan array of object dari studi kasus materi pekan 1--3.}{Menganalisis soal studi kasus, merancang algoritma, menulis program di komputer laboratorium secara mandiri dalam waktu yang ditentukan, dan menjalankan program pada kasus uji.}{Kode program yang terkompilasi beserta hasil eksekusi pada kasus uji, dikumpulkan melalui LMS.}{Program terkompilasi dan berjalan benar pada kasus uji & 8 & Program terkompilasi tanpa error dan menghasilkan keluaran benar untuk seluruh kasus uji yang diberikan. & Program terkompilasi dan berjalan, tetapi keluaran benar hanya pada sebagian kasus uji. & Program tidak terkompilasi atau keluaran salah pada hampir semua kasus uji. \\ \\
Ketepatan implementasi class dan array of object & 7 & Deklarasi class, atribut, method, konstruktor, dan pemanfaatan array of object sesuai spesifikasi soal. & Struktur class dan array of object sebagian besar sesuai, tetapi terdapat kesalahan pada konstruktor, akses atribut, atau indeks. & Struktur class atau array of object tidak sesuai spesifikasi atau tidak digunakan. \\ \\
Kualitas kode dan penjelasan & 5 & Kode terstruktur rapi dengan penamaan konsisten dan mahasiswa mampu menjelaskan alur programnya saat ditanya. & Kode berjalan tetapi kurang terstruktur, atau penjelasan alur program kurang lengkap. & Kode sulit dibaca dan mahasiswa tidak dapat menjelaskan alur programnya. \\}

\assessmentblock{Tes Case Method 1: Algoritma Sorting dan Searching}{Ujian case method (Asessment 2)}{SCPMK0903-01801}{Mahasiswa menyelesaikan kasus pengolahan data dari materi pekan 5--7 dengan memilih dan mengimplementasikan algoritma brute force/divide-conquer, sorting, dan searching yang sesuai.}{Menganalisis kasus, memilih algoritma beserta alasannya, mengimplementasikan penyelesaian dalam flowchart/pseudocode/program, menguji hasil, dan mendokumentasikan langkah penyelesaian secara mandiri.}{Lembar jawaban berisi analisis kasus, flowchart/pseudocode/kode program, hasil pengujian, dan dokumentasi langkah penyelesaian.}{Ketepatan analisis kasus dan pemilihan algoritma & 9 & Karakteristik kasus dianalisis benar dan algoritma sorting/searching yang dipilih tepat disertai alasan berbasis kompleksitas. & Analisis kasus cukup tepat, tetapi alasan pemilihan algoritma lemah atau tidak mempertimbangkan kompleksitas. & Analisis kasus salah atau algoritma yang dipilih tidak sesuai permasalahan. \\ \\
Kebenaran implementasi algoritma pada kasus uji & 10 & Implementasi sorting/searching berjalan benar pada seluruh kasus uji, termasuk simulasi manual langkah per langkah yang tepat. & Implementasi benar pada sebagian kasus uji atau terdapat kesalahan pada beberapa langkah simulasi. & Implementasi salah secara mendasar atau simulasi langkah tidak sesuai algoritma. \\ \\
Kualitas dokumentasi langkah penyelesaian & 6.5 & Langkah penyelesaian didokumentasikan runtut, lengkap dari analisis sampai hasil uji, dan mudah diverifikasi. & Dokumentasi tersedia, tetapi sebagian langkah tidak runtut atau hasil uji tidak lengkap. & Dokumentasi tidak menggambarkan proses penyelesaian atau tidak ada. \\}

\assessmentblock{Tes Case Method 2: Struktur Data Linear}{Ujian case method (Asessment 1)}{SCPMK0708-01801}{Mahasiswa menyelesaikan kasus komputasi dari materi pekan 9--12 dengan memilih dan mengimplementasikan struktur data linear (stack, queue, linked list, double linked list) yang sesuai.}{Menganalisis kebutuhan kasus, memilih struktur data linear yang tepat, mengimplementasikan operasi-operasinya dalam program/flowchart, menguji hasil, dan mendemonstrasikan serta menjelaskan penyelesaian.}{Kode program/flowchart penyelesaian kasus, hasil pengujian, dan demonstrasi/penjelasan singkat penyelesaian.}{Ketepatan pemilihan struktur data linear sesuai kasus & 9 & Struktur data yang dipilih sesuai karakteristik kasus (LIFO/FIFO/dinamis) dengan alasan yang benar. & Struktur data yang dipilih dapat menyelesaikan kasus, tetapi alasan pemilihan kurang tepat atau kurang efisien. & Struktur data yang dipilih tidak sesuai kasus atau tanpa alasan. \\ \\
Kebenaran implementasi operasi struktur data & 10 & Operasi push/pop, enqueue/dequeue, dan insert/delete node diimplementasikan benar untuk seluruh skenario kasus, termasuk kondisi batas. & Sebagian besar operasi benar, tetapi terdapat kesalahan pada kondisi batas (kosong/penuh, node awal/akhir). & Operasi struktur data salah secara mendasar atau program tidak berjalan. \\ \\
Kualitas demo dan penjelasan kode & 6.5 & Mahasiswa mendemonstrasikan program berjalan dan menjelaskan alur setiap operasi struktur data dengan benar. & Demo berjalan, tetapi penjelasan alur operasi kurang lengkap atau sebagian keliru. & Tidak dapat mendemonstrasikan program atau tidak dapat menjelaskan kode yang ditulis. \\}

\assessmentblock{Tes Praktik 2: Struktur Data Non-Linear}{Ujian praktik (Asessment 1)}{SCPMK0703-01801}{Mahasiswa mengerjakan ujian praktik individu di laboratorium untuk mengimplementasikan struktur data non-linear (tree, binary search tree, graph) dari studi kasus materi pekan 14--15.}{Menganalisis soal studi kasus, memilih representasi struktur data non-linear, menulis dan menjalankan program secara mandiri di laboratorium, serta menjelaskan hasil implementasi.}{Kode program yang terkompilasi beserta hasil eksekusi pada kasus uji, dikumpulkan melalui LMS.}{Kebenaran implementasi tree dan binary search tree & 8 & Operasi pembentukan, penyisipan, dan penelusuran tree/BST diimplementasikan benar sesuai tahapan yang tepat. & Implementasi tree/BST sebagian benar, tetapi terdapat kesalahan pada penyisipan atau penelusuran. & Implementasi tree/BST salah secara mendasar atau tidak selesai. \\ \\
Kebenaran implementasi graph & 7 & Representasi graph dan konversi matriks--graph diimplementasikan benar untuk seluruh kasus uji. & Representasi graph benar, tetapi konversi atau penelusuran salah pada sebagian kasus. & Representasi graph salah atau program tidak berjalan. \\ \\
Kebenaran hasil pada kasus uji dan penjelasan kode & 5 & Program menghasilkan keluaran benar pada seluruh kasus uji dan mahasiswa mampu menjelaskan pilihan struktur datanya. & Keluaran benar pada sebagian kasus uji atau penjelasan pilihan struktur data kurang tepat. & Keluaran salah pada hampir semua kasus uji dan mahasiswa tidak dapat menjelaskan kodenya. \\}

\begin{tblr}{width=\textwidth,colspec={|Q[l,m,3.2cm]|X[l,m]|},hlines,vlines,hspan=minimal,row{1-3}={bg=headgray},rowsep=1.5pt,colsep=3pt}
Jadwal Pelaksanaan: & Minggu 1--17 sesuai RPS. Setiap evaluasi memiliki RTM dan rubrik multi-kriteria. \\
Lain-Lain Yang Diperlukan: & LMS, komputer laboratorium, IDE Java/JDK, aplikasi flowchart, dan perangkat presentasi. \\
Pustaka: & Mengacu pustaka utama dan pendukung pada bagian RPS. \\
\end{tblr}
